%% ==================================================
%% Section: Teleportation and long-range gates
%% Sources: arXiv:2307.01696 (Malz, Styliaris, Wei, Cirac), eq. (16) and
%% paragraph "Tree-RG circuit with measurements"; arXiv:2103.13367 (Piroli,
%% Styliaris, Cirac), paragraph "State transformations with QC and LOCC" and
%% Example 2.
%% ==================================================

\section{Teleportation and long-range gates}
\label{sec:teleport}

Registers that are far apart ``can be teleported at neighboring registers with
a constant overhead'' by ``creating nearest-neighbor entangled pairs, then
performing simultaneous measurements, and correcting (without postselection)
based on the measurement outcomes'' \cite{Malz2024Preparation}. This section
proves this for qudits of any dimension: one hop moves the content of a site
across two sites through an entangled pair on two cells of any geometry, any
number of sites is moved at once along chains of hops of any length in depth at
most $2$, and on the ring two-site gates joining the sites $a$ and $a+2L+1$, on
pairwise disjoint stretches $\{a,\dots,a+2L+1\}$, are applied in depth at most
$5$ for every $L$. Gates at even separation, or on overlapping stretches, are
not covered. Registers of $s\ge1$ consecutive sites are teleported site by
site, and a layer of gates on pairs of distant registers is applied in depth at
most $4s+K$ when the gates on the registers, brought next to each other, form a
unitary circuit of depth at most $K$. The ancilla teleported through entangled
pairs between neighbouring sites by local operations and classical
communication in \cite{Piroli2021LOCC} is the case of one hop.

Throughout, the local labels are $\mathbb Z_d=\{0,\dots,d-1\}$ with $d\ge1$,
$\zeta=e^{2\pi i/d}$, and $F_{ab}=d^{-1/2}\zeta^{ab}$ is the Fourier matrix.

\begin{definition}[Vectors with zero at given sites]
    \label{def:teleport_zero_on}
    \lean{QuantumCircuit.IsZeroOn, QuantumCircuit.IsZeroOn.mono,
        QuantumCircuit.IsZeroOn.smul, QuantumCircuit.zeroOnSubmodule,
        QuantumCircuit.mem_zeroOnSubmodule}
    \leanok
    A vector $v$ on the configurations of qudits on the sites $S$ has $\ket0$ on
    $D\subseteq S$ if $v(x)=0$ whenever $x_i\ne0$ for some $i\in D$. These
    vectors form a linear subspace, which shrinks as $D$ grows.
\end{definition}

\begin{lemma}[Projections, one-site operators and permutations of sites]
    \label{lem:teleport_site_operators}
    \lean{QuantumCircuit.place_single_self_mulVec_apply,
        QuantumCircuit.place_single_union, QuantumCircuit.commute_place_single,
        QuantumCircuit.prodOp_update_one_mulVec_apply,
        QuantumCircuit.isZeroOn_iff_place_single_mulVec,
        QuantumCircuit.IsZeroOn.mulVec_of_mem_supportedOperators,
        QuantumCircuit.IsZeroOn.prodOp_update_one_mulVec,
        QuantumCircuit.IsZeroOn.sitePermOp_mulVec_image,
        QuantumCircuit.IsZeroOn.sitePermOp_mulVec_of_mapsTo,
        QuantumCircuit.sitePermOp_mul_prodOp}
    \leanok
    \uses{def:teleport_zero_on, def:sites_place, def:hastings_site_permutation}
    Let $P_{D,m}$ be the matrix unit $\ket m\bra m$ of a configuration $m$ of
    $D$, placed on $D$, and let $P_\pi$ act by $(P_\pi v)(x)=v(x\circ\pi)$.
    \begin{enumerate}
        \item $(P_{D,m}v)(x)=v(x)$ if $x$ restricts to $m$ on $D$, and $0$
              otherwise; for the restrictions of one configuration $z$,
              $P_{D\cup D',z}=P_{D,z}P_{D',z}$, and these projections commute.
        \item The operator $u$ at the site $t$ acts by
              $(u_tv)(x)=\sum_ju_{x_t,j}\,v(x[t\mapsto j])$.
        \item A vector has $\ket0$ on $D$ if and only if $P_{D,0}$ fixes it;
              hence every operator supported on a region disjoint from $D$, in
              particular a one-site operator at a site outside $D$, keeps
              $\ket0$ on $D$.
        \item If $v$ has $\ket0$ on $D$, then $P_\pi v$ has $\ket0$ on $\pi(D)$,
              and on $D'$ whenever $\pi$ maps no site outside $D$ into $D'$.
        \item $P_\pi\bigotimes_ig_i=\bigl(\bigotimes_ig_{\pi^{-1}(i)}\bigr)P_\pi$.
    \end{enumerate}
\end{lemma}
\begin{proof}
    \leanok
    Items (1), (2), (4) and (5) follow by evaluating at a configuration. For
    (3), $P_{D,0}$ is supported on $D$, so it commutes with every operator
    supported on a disjoint region.
\end{proof}

\begin{definition}[A teleportation hop]
    \label{def:teleport_hop}
    \lean{QuantumCircuit.TeleportHop, QuantumCircuit.TeleportHop.pairGate₁,
        QuantumCircuit.TeleportHop.pairGate₂,
        QuantumCircuit.TeleportHop.pairGate₁_mem_unitaryGroup,
        QuantumCircuit.TeleportHop.pairGate₂_mem_unitaryGroup,
        QuantumCircuit.TeleportHop.gate₁, QuantumCircuit.TeleportHop.gate₂,
        QuantumCircuit.TeleportHop.gate₁_mem_unitaryGroup,
        QuantumCircuit.TeleportHop.gate₂_mem_unitaryGroup,
        QuantumCircuit.TeleportHop.gate₁_mem_supportedOperators,
        QuantumCircuit.TeleportHop.gate₂_mem_supportedOperators,
        QuantumCircuit.TeleportHop.pre, QuantumCircuit.TeleportHop.frame,
        QuantumCircuit.TeleportHop.frame_mem_unitaryGroup,
        QuantumCircuit.TeleportHop.frame_of_ne}
    \leanok
    \uses{def:circuits_geometry, def:sites_place, def:teleport_zero_on}
    A hop on a geometry consists of three distinct sites $c,e,f$ and two cells
    whose regions are $\{e,f\}$ and $\{c,e\}$. Its first gate $G_1$, on
    $\{e,f\}$, applies $F$ at $e$ and then $\ket a_e\ket b_f\mapsto\ket
        a_e\ket{b+a}_f$; its second gate $G_2$, on $\{c,e\}$, applies $\ket
        a_c\ket b_e\mapsto\ket a_c\ket{b-a}_e$ and then $F$ at $c$. Both are
    unitary. For a configuration $z$, the hop before the correction is
    $P_{\{c,e\},z}G_2G_1$, and its frame consists of the one-site unitaries
    $\ket q\mapsto\ket{q+z_c}$ at $c$, $\ket q\mapsto\ket{q+z_e}$ at $e$,
    $X^{z_e}Z^{z_c}\colon\ket q\mapsto\zeta^{z_cq}\ket{q+z_e}$ at $f$, and the
    identity elsewhere.
\end{definition}

\begin{lemma}[One hop]
    \label{lem:teleport_hop}
    \lean{QuantumCircuit.TeleportHop.pre_mulVec}
    \leanok
    \uses{def:teleport_hop, lem:teleport_site_operators}
    On every vector $v$ with $\ket0$ on $\{e,f\}$,
    \[
        P_{\{c,e\},z}G_2G_1v=\frac1d\Bigl(\bigotimes_i\Phi_i\Bigr)P_{(c\,f)}v,
    \]
    where $\Phi$ is the frame of the hop for $z$ and $(c\,f)$ exchanges $c$
    and $f$.
\end{lemma}
\begin{proof}
    \leanok
    On $\ket0_e\ket0_f$, the gate $G_1$ prepares
    $d^{-1/2}\sum_j\ket j_e\ket j_f$, so
    $(G_1v)(w)=d^{-1/2}v(w[e\mapsto0,f\mapsto w_f-w_e])$. Applying $G_2$ gives
    $(G_2G_1v)(y)=F_{y_c,\,y_f-y_e}\,d^{-1/2}\,v(y[c\mapsto y_f-y_e,e\mapsto0,f\mapsto0])$.
    For $(y_c,y_e)=(z_c,z_e)$ the amplitude is
    $d^{-1}\zeta^{z_c(y_f-z_e)}$, the value read at $c$ is the label $y_f-z_e$
    carried to $f$, and the right-hand side evaluates to the same; otherwise
    both sides vanish, the right-hand side because $v$ is read at a
    configuration with a nonzero label at $e$ or $f$.
\end{proof}

\begin{definition}[Layers given by lists]
    \label{def:teleport_list_layer}
    \lean{QuantumCircuit.Layer.ofList, QuantumCircuit.Layer.kraus_ofList,
        QuantumCircuit.Layer.forall_isUnitary_ofList}
    \leanok
    \uses{def:circuits_of_cells, thm:circuits_of_cells, def:sites_place}
    For a list of operators $G_i$ supported on pairwise disjoint nonempty
    cells, the layer with, on each of these cells, the matrix whose placement is
    the corresponding $G_i$, and the identity elsewhere, has the Kraus operator
    $\prod_iG_i$, the product along the list. It is unitary when the $G_i$ are.
\end{definition}

\begin{definition}[Valid lists of hops]
    \label{def:teleport_valid}
    \lean{QuantumCircuit.TeleportHop.sites,
        QuantumCircuit.TeleportHop.allSites,
        QuantumCircuit.TeleportHop.pairSites,
        QuantumCircuit.TeleportHop.measuredSites,
        QuantumCircuit.TeleportHop.Valid, QuantumCircuit.TeleportHop.valid_nil,
        QuantumCircuit.TeleportHop.mem_allSites,
        QuantumCircuit.TeleportHop.measuredSites_subset_allSites,
        QuantumCircuit.TeleportHop.disjoint_measuredSites,
        QuantumCircuit.TeleportHop.pairwise_disjoint_one,
        QuantumCircuit.TeleportHop.pairwise_disjoint_two,
        QuantumCircuit.TeleportHop.layerOne, QuantumCircuit.TeleportHop.layerTwo,
        QuantumCircuit.TeleportHop.forall_isUnitary_layerOne,
        QuantumCircuit.TeleportHop.forall_isUnitary_layerTwo,
        QuantumCircuit.TeleportHop.kraus_layerOne,
        QuantumCircuit.TeleportHop.kraus_layerTwo,
        QuantumCircuit.TeleportHop.sitePerm,
        QuantumCircuit.TeleportHop.sitePerm_apply_of_notMem,
        QuantumCircuit.IsZeroOn.sitePermOp_sitePerm_mulVec}
    \leanok
    \uses{def:teleport_hop, def:teleport_list_layer}
    A list of hops $(h_1,\dots,h_H)$, the most recent first, is valid if for
    every $j<k$ the sites $e_j,f_j$ are not sites of $h_k$ and $c_j$ is neither
    $c_k$ nor $e_k$; the site $c_j$ may be $f_k$, so hops chain as
    $c_k\to f_k=c_j\to f_j$. For a valid list the cells $\{e_j,f_j\}$ are
    pairwise disjoint, and so are the cells $\{c_j,e_j\}$: the first gates of
    all hops form one unitary layer, and so do the second gates. The
    permutation of sites of the list is $(c_1\,f_1)\cdots(c_H\,f_H)$; it fixes
    every site of no hop, so $P_\pi$ keeps $\ket0$ on any set of sites met by
    no hop.
\end{definition}

\begin{lemma}[Chains of hops before the correction]
    \label{lem:teleport_chain_pre}
    \lean{QuantumCircuit.TeleportHop.chainFrame,
        QuantumCircuit.TeleportHop.chainPre,
        QuantumCircuit.TeleportHop.chainFrame_mem_unitaryGroup,
        QuantumCircuit.TeleportHop.chainFrame_of_notMem,
        QuantumCircuit.TeleportHop.prod_gate₂_mem_supportedOperators,
        QuantumCircuit.TeleportHop.pre_mem_supportedOperators,
        QuantumCircuit.TeleportHop.chainPre_cons,
        QuantumCircuit.TeleportHop.chainPre_mulVec}
    \leanok
    \uses{def:teleport_valid, lem:teleport_hop, lem:teleport_site_operators}
    For a valid list of $H$ hops and a configuration $z$, there are one-site
    unitaries $\Phi_i$, equal to the identity off the sites of the hops, such
    that on every vector $v$ with $\ket0$ at the sites $e_j,f_j$ of all hops
    \[
        P_{M,z}\Bigl(\prod_jG_{2,h_j}\Bigr)\Bigl(\prod_jG_{1,h_j}\Bigr)v
        =d^{-H}\Bigl(\bigotimes_i\Phi_i\Bigr)P_\pi v,
    \]
    $M$ the set of the sites $c_j,e_j$ and $\pi$ the permutation of sites of
    the list.
\end{lemma}
\begin{proof}
    \leanok
    By induction on the list. The first gate of $h_1$ commutes with the second
    gates of the later hops, whose sites it does not meet, and the projection
    onto the later outcomes commutes with the operator of $h_1$, so the
    operator of the list is that of $h_1$ times that of the rest. The frame of
    the rest is split into its unitary at $c_1$ and the others; the others
    commute with the operator of $h_1$, the unitary at $c_1$ keeps $\ket0$ on
    $\{e_1,f_1\}$, and Lemma~\ref{lem:teleport_hop} moves it to $f_1$ through
    the exchange $(c_1\,f_1)$.
\end{proof}

\begin{theorem}[Teleportation in one round]
    \label{thm:teleport_round}
    \lean{QuantumCircuit.TeleportHop.roundUnitary,
        QuantumCircuit.TeleportHop.extendOutcome,
        QuantumCircuit.TeleportHop.round,
        QuantumCircuit.TeleportHop.op_roundUnitary,
        QuantumCircuit.TeleportHop.depth_round_le,
        QuantumCircuit.TeleportHop.depth_round_id_le,
        QuantumCircuit.TeleportHop.forall_round,
        QuantumCircuit.TeleportHop.isImplementationOn_round}
    \leanok
    \uses{lem:teleport_chain_pre, def:measurement_circuit,
        def:measurement_correct, lem:measurement_round_implementation,
        def:measurement_implementation}
    For a valid list of $H$ hops and a unitary circuit $u$ with operator $U$,
    the teleportation round applies $u$, the layer of the first gates and the
    layer of the second gates of the hops, measures the sites $c_j,e_j$, and
    corrects by the inverse of the frame of Lemma~\ref{lem:teleport_chain_pre}
    for the outcome. It is a measurement circuit of depth at most
    $\operatorname{depth}(u)+2$, whatever the number and the lengths of the
    chains, and it implements $P_\pi U$ on the vectors $v$ such that $Uv$ has
    $\ket0$ at the sites $e_j,f_j$ of all hops, with the scalar $d^{-H}$ on
    every branch.
\end{theorem}
\begin{proof}
    \leanok
    The unitary part of the round has the operator
    $\bigl(\prod_jG_{2,h_j}\bigr)\bigl(\prod_jG_{1,h_j}\bigr)U$. For the
    outcome $m$, extended by $0$ to a configuration $z$, the measurement
    projection is $P_{M,z}$, and by Lemma~\ref{lem:teleport_chain_pre} the
    corrected operator maps $v$ to
    $(\bigotimes_i\Phi_i)^\dagger d^{-H}(\bigotimes_i\Phi_i)P_\pi Uv=d^{-H}P_\pi
        Uv$. Lemma~\ref{lem:measurement_round_implementation} concludes.
\end{proof}

\begin{lemma}[Gates conjugated by teleportation]
    \label{lem:teleport_conjugated}
    \lean{QuantumCircuit.TeleportHop.isImplementationOn_round_comp_round}
    \leanok
    \uses{thm:teleport_round, lem:measurement_implementation_comp}
    For valid lists of hops $\mathcal T$ and $\mathcal B$ and a unitary circuit
    with operator $U$, the round teleporting along $\mathcal T$ followed by
    the round applying the circuit and teleporting along $\mathcal B$
    implements $P_{\mathcal B}UP_{\mathcal T}$ on the vectors $v$ with $\ket0$
    at the sites $e,f$ of $\mathcal T$ such that $UP_{\mathcal T}v$ has $\ket0$
    at the sites $e,f$ of $\mathcal B$.
\end{lemma}
\begin{proof}
    \leanok
    Theorem~\ref{thm:teleport_round} for both rounds, the first after the empty
    circuit, composed by Lemma~\ref{lem:measurement_implementation_comp}.
\end{proof}

\begin{lemma}[Chains between distant sites of the ring]
    \label{lem:teleport_chains}
    \lean{QuantumCircuit.add_natCast_injective, QuantumCircuit.add_natCast_succ,
        QuantumCircuit.add_natCast_ne,
        QuantumCircuit.TeleportHop.After,
        QuantumCircuit.TeleportHop.valid_iff_pairwise,
        QuantumCircuit.TeleportHop.sitePerm_append,
        QuantumCircuit.TeleportHop.pairSites_append,
        QuantumCircuit.TeleportHop.chainHop,
        QuantumCircuit.TeleportHop.chainHopBack,
        QuantumCircuit.TeleportHop.forwardChain,
        QuantumCircuit.TeleportHop.backwardChain,
        QuantumCircuit.TeleportHop.mem_forwardChain,
        QuantumCircuit.TeleportHop.mem_backwardChain,
        QuantumCircuit.TeleportHop.valid_forwardChain,
        QuantumCircuit.TeleportHop.valid_backwardChain,
        QuantumCircuit.TeleportHop.sitePerm_forwardChain,
        QuantumCircuit.TeleportHop.sitePerm_backwardChain,
        QuantumCircuit.TeleportHop.pairSites_forwardChain_subset,
        QuantumCircuit.TeleportHop.pairSites_backwardChain_subset,
        QuantumCircuit.TeleportHop.allSites_forwardChain_subset,
        QuantumCircuit.TeleportHop.allSites_backwardChain_subset,
        QuantumCircuit.TeleportHop.isZeroOn_sitePermOp_forwardChain}
    \leanok
    \uses{def:teleport_valid, def:circuits_ring_geometry,
        lem:teleport_site_operators}
    On the ring of $N$ sites, with cells $\{k,k+1\}$, let $2L<N$. The forward
    chain from $a$ consists of the hops $(a+2j,a+2j+1,a+2j+2)$, $j=L-1,\dots,0$,
    and the backward chain of the hops $(a+2j+2,a+2j+1,a+2j)$,
    $j=0,\dots,L-1$. Both are valid, their sites lie in $\{a,\dots,a+2L\}$, and
    their permutations of sites are inverse to each other; the forward one maps
    $a$ to $a+2L$ and fixes $a+j$ for $2L<j<N$. The sites $e,f$ of the forward
    chain lie in $\{a+1,\dots,a+2L\}$ and those of the backward chain in
    $\{a,\dots,a+2L-1\}$, and if $v$ has $\ket0$ on $\{a+1,\dots,a+2L\}$, the
    forward chain moves it to a vector with $\ket0$ on $\{a,\dots,a+2L-1\}$.
\end{lemma}
\begin{proof}
    \leanok
    The sites $a+j$, $0\le j<N$, are distinct. The transpositions
    $(a+2j\;\;a+2j+2)$, $j=0,\dots,L-1$, move $a$ to $a+2L$ and shift
    $\{a+1,\dots,a+2L\}$ onto $\{a,\dots,a+2L-1\}$, and the backward chain
    applies them in the reverse order.
\end{proof}

\begin{lemma}[Parallel chains]
    \label{lem:teleport_parallel}
    \lean{QuantumCircuit.TeleportHop.mem_allSites_flatMap,
        QuantumCircuit.TeleportHop.mem_pairSites_flatMap,
        QuantumCircuit.TeleportHop.after_of_disjoint,
        QuantumCircuit.TeleportHop.sites_subset_allSites,
        QuantumCircuit.TeleportHop.valid_flatMap,
        QuantumCircuit.TeleportHop.sitePerm_flatMap_apply,
        QuantumCircuit.TeleportHop.sitePerm_flatMap_apply_of_notMem,
        Equiv.Perm.apply_mem_of_forall_notMem}
    \leanok
    \uses{def:teleport_valid, lem:teleport_chains}
    If valid lists of hops $\mathcal L_g$ have their sites in pairwise disjoint
    sets $S_g$, their concatenation is valid, and its permutation of sites
    agrees with that of $\mathcal L_g$ on $S_g$ and fixes every site outside
    $\bigcup_gS_g$.
\end{lemma}
\begin{proof}
    \leanok
    Hops on disjoint sets of sites satisfy the validity condition in either
    order, and the permutation of sites of each $\mathcal L_g$ maps $S_g$ into
    itself and fixes the other sites.
\end{proof}

\begin{theorem}[A layer of distant gates in depth five with measurements]
    \label{thm:teleport_long_range}
    \lean{QuantumCircuit.LongRangeGate, QuantumCircuit.LongRangeGate.far,
        QuantumCircuit.LongRangeGate.near, QuantumCircuit.LongRangeGate.span,
        QuantumCircuit.LongRangeGate.interior,
        QuantumCircuit.LongRangeGate.cleared, QuantumCircuit.LongRangeGate.op,
        QuantumCircuit.LongRangeGate.mem_span,
        QuantumCircuit.LongRangeGate.mem_interior,
        QuantumCircuit.LongRangeGate.localGate,
        QuantumCircuit.LongRangeGate.there, QuantumCircuit.LongRangeGate.back,
        QuantumCircuit.LongRangeGate.near_ne_far,
        QuantumCircuit.LongRangeGate.a_ne_far,
        QuantumCircuit.LongRangeGate.near_add_one,
        QuantumCircuit.LongRangeGate.localGate_mem_unitaryGroup,
        QuantumCircuit.LongRangeGate.localGate_mem_supportedOperators,
        QuantumCircuit.LongRangeGate.op_mem_unitaryGroup,
        QuantumCircuit.LongRangeGate.op_mem_supportedOperators,
        QuantumCircuit.LongRangeGate.far_mem_span,
        QuantumCircuit.LongRangeGate.cell_near_subset_span,
        QuantumCircuit.LongRangeGate.interior_subset_span,
        QuantumCircuit.LongRangeGate.cleared_subset_span,
        QuantumCircuit.LongRangeGate.allSites_there_subset_span,
        QuantumCircuit.LongRangeGate.allSites_back_subset_span,
        QuantumCircuit.LongRangeGate.disjoint_cleared_cell_near,
        QuantumCircuit.LongRangeGate.sitePerm_back_near,
        QuantumCircuit.LongRangeGate.sitePerm_back_far,
        QuantumCircuit.LongRangeGate.sitePerm_back_sitePerm_there,
        QuantumCircuit.LongRangeGate.pairwise_disjoint_cell,
        QuantumCircuit.LongRangeGate.localLayer,
        QuantumCircuit.LongRangeGate.kraus_localLayer,
        QuantumCircuit.LongRangeGate.forall_isUnitary_localLayer,
        QuantumCircuit.LongRangeGate.valid_there,
        QuantumCircuit.LongRangeGate.valid_back,
        QuantumCircuit.LongRangeGate.isZeroOn_sitePermOp_there,
        QuantumCircuit.LongRangeGate.isZeroOn_localGates,
        QuantumCircuit.LongRangeGate.sitePerm_there_eq_symm,
        QuantumCircuit.LongRangeGate.sitePermOp_back_mul_localGates_mul_sitePermOp_there,
        QuantumCircuit.LongRangeGate.rounds,
        QuantumCircuit.LongRangeGate.depth_rounds_le,
        QuantumCircuit.LongRangeGate.forall_rounds,
        QuantumCircuit.LongRangeGate.isImplementationOn_rounds,
        List.mul_prod_mul_of_mul_eq_one}
    \leanok
    \uses{lem:teleport_conjugated, lem:teleport_chains, lem:teleport_parallel,
        def:sites_place_along, lem:sites_place_along_permutation, def:teleport_list_layer}
    On the ring of $N$ sites, let $u_1,\dots,u_n$ be two-site unitaries on the
    pairs $\{a_r,a_r+2L_r+1\}$, $2L_r+1<N$, at odd separation, whose stretches
    $\{a_r,\dots,a_r+2L_r+1\}$ are pairwise disjoint. A measurement circuit of
    depth at most $5$ implements $\prod_ru_r$ on the vectors with $\ket0$ at
    the sites $a_r+1,\dots,a_r+2L_r$ of every $r$, without postselection.
\end{theorem}
\begin{proof}
    \leanok
    By Lemmas~\ref{lem:teleport_chains} and~\ref{lem:teleport_parallel}, the
    concatenated forward chains $\mathcal T$ from the sites $a_r$ and the
    concatenated backward chains $\mathcal B$ are valid with
    $P_{\mathcal B}=P_{\mathcal T}^{-1}$, and $P_{\mathcal T}v$ has $\ket0$ at
    the sites $a_r,\dots,a_r+2L_r-1$, which the layer
    $G=\prod_r(u_r)_{\{a_r+2L_r,\,a_r+2L_r+1\}}$ on neighbouring pairs keeps.
    By Lemma~\ref{lem:teleport_conjugated}, the round of $\mathcal T$ followed
    by the round of $G$ and $\mathcal B$, of depths at most $2$ and $3$,
    implements $P_{\mathcal T}^{-1}GP_{\mathcal T}$, which is
    $\prod_r(u_r)_{\{a_r,\,a_r+2L_r+1\}}$ by
    Lemma~\ref{lem:sites_place_along_permutation}, since $P_{\mathcal B}$ maps
    $a_r+2L_r$ to $a_r$ and fixes $a_r+2L_r+1$.
\end{proof}

\begin{definition}[Register gates]
    \label{def:teleport_register_gate}
    \lean{QuantumCircuit.RegisterGate, QuantumCircuit.RegisterGate.offset,
        QuantumCircuit.RegisterGate.sites,
        QuantumCircuit.RegisterGate.localSites,
        QuantumCircuit.RegisterGate.span, QuantumCircuit.RegisterGate.zone,
        QuantumCircuit.RegisterGate.interior,
        QuantumCircuit.RegisterGate.mem_span,
        QuantumCircuit.RegisterGate.mem_zone,
        QuantumCircuit.RegisterGate.offset_lt,
        QuantumCircuit.RegisterGate.sites_injective,
        QuantumCircuit.RegisterGate.localSites_injective,
        QuantumCircuit.RegisterGate.localSites_succ,
        QuantumCircuit.RegisterGate.sitesEmb,
        QuantumCircuit.RegisterGate.localSitesEmb,
        QuantumCircuit.RegisterGate.op, QuantumCircuit.RegisterGate.localOp,
        QuantumCircuit.RegisterGate.sites_mem_span,
        QuantumCircuit.RegisterGate.localSites_mem_span,
        QuantumCircuit.RegisterGate.zone_subset_span,
        QuantumCircuit.RegisterGate.op_mem_unitaryGroup,
        QuantumCircuit.RegisterGate.op_mem_supportedOperators,
        QuantumCircuit.RegisterGate.add_eq_add_iff}
    \leanok
    \uses{def:sites_place_along, thm:sites_place_along_support, def:circuits_ring_geometry}
    On the ring of $N$ sites, a register gate $(X,a,L)$ with $2L+2s\le N$
    places a matrix $X$ on $2s$ qudits along the two registers
    $\{a,\dots,a+s-1\}$ and $\{a+2L+s,\dots,a+2L+2s-1\}$; it is unitary when
    $X$ is. Its stretch is $\{a,\dots,a+2L+2s-1\}$, its zones are
    $Z_t=\{a+t,\dots,a+t+2L-1\}$, $0\le t\le s$, its interior is $Z_s$, and its
    local gate places $X$ along the consecutive sites
    $a+2L,\dots,a+2L+2s-1$.
\end{definition}

\begin{lemma}[Teleporting registers site by site]
    \label{lem:teleport_register_chains}
    \lean{QuantumCircuit.RegisterGate.two_mul_L_lt,
        QuantumCircuit.RegisterGate.there, QuantumCircuit.RegisterGate.back,
        QuantumCircuit.RegisterGate.allSites_there_subset,
        QuantumCircuit.RegisterGate.allSites_back_subset,
        QuantumCircuit.RegisterGate.allSites_there_subset_span,
        QuantumCircuit.RegisterGate.allSites_back_subset_span,
        QuantumCircuit.RegisterGate.pairSites_there_subset,
        QuantumCircuit.RegisterGate.pairSites_back_subset,
        QuantumCircuit.RegisterGate.sitePerm_there_apply_of_lt,
        QuantumCircuit.RegisterGate.sitePerm_there_apply_self,
        QuantumCircuit.RegisterGate.sitePerm_back_eq_inv,
        QuantumCircuit.RegisterGate.sitePerm_back_apply_of_lt,
        QuantumCircuit.RegisterGate.sitePerm_back_apply_self,
        QuantumCircuit.RegisterGate.sitePerm_there_notMem_zone,
        QuantumCircuit.RegisterGate.sitePerm_back_notMem_zone,
        QuantumCircuit.RegisterGate.thereAll,
        QuantumCircuit.RegisterGate.backAll,
        QuantumCircuit.RegisterGate.zoneAll,
        QuantumCircuit.RegisterGate.mem_zoneAll,
        QuantumCircuit.RegisterGate.valid_thereAll,
        QuantumCircuit.RegisterGate.valid_backAll,
        QuantumCircuit.RegisterGate.sitePerm_thereAll_apply,
        QuantumCircuit.RegisterGate.sitePerm_backAll_apply,
        QuantumCircuit.RegisterGate.isZeroOn_sitePermOp_thereAll,
        QuantumCircuit.RegisterGate.isZeroOn_sitePermOp_backAll,
        QuantumCircuit.RegisterGate.pairSites_thereAll_subset,
        QuantumCircuit.RegisterGate.pairSites_backAll_subset,
        QuantumCircuit.RegisterGate.sitePerm_backAll_eq_inv,
        QuantumCircuit.RegisterGate.therePerm,
        QuantumCircuit.RegisterGate.therePerm_apply,
        QuantumCircuit.RegisterGate.therePerm_comp_sites,
        QuantumCircuit.RegisterGate.isZeroOn_sitePermOp_therePerm,
        QuantumCircuit.RegisterGate.isZeroOn_sitePermOp_therePerm_inv,
        QuantumCircuit.RegisterGate.thereRounds,
        QuantumCircuit.RegisterGate.backRounds,
        QuantumCircuit.RegisterGate.depth_thereRounds_le,
        QuantumCircuit.RegisterGate.depth_backRounds_le,
        QuantumCircuit.RegisterGate.forall_thereRounds,
        QuantumCircuit.RegisterGate.forall_backRounds,
        QuantumCircuit.RegisterGate.isImplementationOn_thereRounds,
        QuantumCircuit.RegisterGate.isImplementationOn_backRounds}
    \leanok
    \uses{def:teleport_register_gate, thm:teleport_round, lem:teleport_chains,
        lem:teleport_parallel, lem:teleport_site_operators,
        lem:measurement_implementation_comp}
    Let $g_1,\dots,g_n$ be register gates with pairwise disjoint stretches and
    $s\ge1$. For $t=s-1,\dots,0$, the round of the forward chains from the
    sites $a_r+t$ teleports $a_r+t$ to $a_r+2L_r+t$ and moves $\ket0$ from
    $\bigcup_rZ_{r,t+1}$ to $\bigcup_rZ_{r,t}$. The $s$ forward rounds form a
    measurement circuit of depth at most $2s$ which implements the composite
    $P$ of their permutations of sites on the vectors with $\ket0$ on
    $\bigcup_rZ_{r,s}$, and $P$ moves these vectors to vectors with $\ket0$ on
    $\bigcup_rZ_{r,0}$; the $s$ backward rounds, for $t=0,\dots,s-1$, form a
    measurement circuit of depth at most $2s$ which implements $P^{-1}$ on the
    vectors with $\ket0$ on $\bigcup_rZ_{r,0}$. The composite $P$ maps the
    sites of the two registers of $g_r$ to its consecutive sites.
\end{lemma}
\begin{proof}
    \leanok
    By Lemmas~\ref{lem:teleport_chains} and~\ref{lem:teleport_parallel}, each
    round is valid, requires $\ket0$ on $\bigcup_rZ_{r,t+1}$, and maps no site
    outside $\bigcup_rZ_{r,t+1}$ into $\bigcup_rZ_{r,t}$. The rounds compose by
    Theorem~\ref{thm:teleport_round} and
    Lemma~\ref{lem:measurement_implementation_comp}; by induction on $t$, the
    forward rounds move $a_r+u$ to $a_r+2L_r+u$ for $u<t$ and fix $a_r+j$ for
    $j\ge t+2L_r$ in the stretch.
\end{proof}

\begin{theorem}[A layer of register gates in constant depth with measurements]
    \label{thm:teleport_register_layer}
    \lean{QuantumCircuit.RegisterGate.rounds,
        QuantumCircuit.RegisterGate.depth_rounds_le,
        QuantumCircuit.RegisterGate.forall_rounds,
        QuantumCircuit.RegisterGate.localProd,
        QuantumCircuit.RegisterGate.isZeroOn_localProd_mulVec,
        QuantumCircuit.RegisterGate.sitePermOp_inv_mul_localProd_mul_sitePermOp,
        QuantumCircuit.RegisterGate.isZeroOn_zoneAll_iff,
        QuantumCircuit.RegisterGate.isImplementationOn_rounds,
        QuantumCircuit.RegisterGate.isUnitaryCircuit_localProd,
        QuantumCircuit.RegisterGate.exists_rounds}
    \leanok
    \uses{lem:teleport_register_chains, lem:measurement_implementation_comp,
        def:measurement_circuit, lem:sites_place_along_permutation}
    Let $g_1,\dots,g_n$ be register gates with pairwise disjoint stretches, and
    let the product of their local gates be the operator of a unitary circuit
    of depth at most $K$ on the ring. Then a measurement circuit of depth at
    most $4s+K$ implements $\prod_rg_r$ on the vectors with $\ket0$ on the
    interiors of all $g_r$, without postselection. Such a circuit exists when
    every $X_r$ is a product of gates of total cost at most $K$ on neighbouring
    sites of the open chain of $2s$ sites.
\end{theorem}
\begin{proof}
    \leanok
    The local gates are supported outside $\bigcup_rZ_{r,0}$ and keep $\ket0$
    there. The $s$ forward rounds, the unitary circuit and the $s$ backward
    rounds of Lemma~\ref{lem:teleport_register_chains} compose by
    Lemma~\ref{lem:measurement_implementation_comp} to an implementation of
    $P^{-1}\bigl(\prod_r\mathrm{place}_{\mathrm{local}}(X_r)\bigr)P$, which is
    $\prod_rg_r$ by Lemma~\ref{lem:sites_place_along_permutation}, in depth at
    most $2s+K+2s$. For the last claim, each product of gates on the open chain
    placed on the consecutive sites of $g_r$ is a unitary circuit of depth at
    most $K$ inside them, and circuits inside pairwise disjoint regions run in
    the same layers.
\end{proof}
